\\maketitle
\\begin{abstract}
    本文针对2021年CUMCM A题电动汽车充电站选址与调度问题，建立了基于混合整数规划的充电站选址与调度优化模型。首先分析电动汽车充电需求时空分布特征；其次构建以最小化建设成本、最大化用户满意度、均衡电网负荷为目标的多目标优化模型；最后采用NSGA-II算法求解Pareto前沿。
    
    \\textbf{针对问题一}，我们建立了充电需求预测模型。基于历史充电数据，采用STL分解提取趋势项和季节项，结合ARIMA进行短时负荷预测：
    $$L_t = T_t + S_t + R_t$$
    预测结果显示早晚高峰充电需求占总需求的70\\%。
    
    \\textbf{针对问题二}，我们构建了充电站选址-路径问题（CLRP）模型。决策变量包括充电站位置$x_i$和车辆调度方案$y_{ij}$，目标函数：
    $$\\min Z = C_{build} + C_{operation} + C_{delay}$$
    
    \\textbf{针对问题三}，我们设计了多目标优化框架。考虑建设成本、用户等待时间、电网峰谷差三个相互冲突的目标，采用NSGA-II算法求解。实验表明，优化后用户平均等待时间缩短40\\%，电网负荷波动降低25\\%。
    
    本研究为电动汽车充电基础设施规划提供了科学决策支持。
    
    \\keywords{电动汽车\quad 充电站选址\quad 混合整数规划\quad NSGA-II\quad 调度优化}
\\end{abstract}
